
AlphaFold2 made accurate protein structure prediction routine
\citep{jumper2021alphafold}, and in doing so turned it into a computing
workload: large screening campaigns run inference on thousands of
sequences, and each prediction is a deep network evaluated on
specialized hardware. Its public implementation is written in JAX, which
compiles the same Python program through XLA for CPUs, GPUs and Google's
Tensor Processing Units (TPUs). On paper this makes the choice of
accelerator a matter of picking the fastest or cheapest device. In
practice, the performance a user obtains depends on more than the
device: on how the program is compiled, on how its computation is mapped
onto the chips that are allocated, and on how stable the underlying
cloud machine is from one session to the next. For this class of
scientific workload, the gap between what the hardware can deliver and
what a user running the model actually gets has received little
systematic attention.

This paper measures that gap. We run the same AlphaFold2 inference code
on the same input on three backends: a Colab CPU runtime, an NVIDIA T4
GPU, and a dedicated eight-chip Google Cloud TPU~v5e slice. We then ask
three questions: how the backends compare on a single call, where the
time goes when a call is first compiled, and what it takes to put all
the allocated chips to work.

The single-call comparison favors the TPU by a wide margin: in steady
state, one TPU chip completes an inference call in 0.47\,s, 27.8$\times$
faster than the T4 in our August baseline
(Section~\ref{sec:results-hardware}).
That comparison comes with three qualifications. It is the performance of
\emph{one} chip: in the default execution path, seven of the eight
allocated chips hold no data at all, so a user pays for a full slice and
uses an eighth of it, and at list prices this idle capacity makes the
TPU slice cost about the same per prediction as the T4 GPU on that same
August baseline (Section~\ref{sec:cost}). The first call on each new
input shape is dominated by work that precedes execution: in a profiler
trace of the first call, about three quarters of the traced
model-application span is self time in JAX's tracing and compilation
path (Section~\ref{sec:profiling}). For one-off jobs, this cold start
matters more than steady-state speed. Finally, the CPU and GPU
baselines did not reproduce across Colab sessions weeks apart: the CPU became roughly
1.6--1.7$\times$ slower and the GPU about twice as fast, for reasons we
could not identify. The GPU change bears directly on the TPU-over-GPU
ratio above, so we report every session rather than a single canonical
value (Section~\ref{sec:rh-variability}).

The idle chips are not recovered by the obvious fix. Vectorizing the
model over a batch of queries with \texttt{jax.vmap} does not help:
throughput never exceeds that of a single query and falls to roughly
three quarters of it at batch sizes 4 and 8 (0.73$\times$ and
0.74$\times$). Mapping independent queries
across devices with \texttt{jax.pmap} does put all eight chips to work:
on a matched grid, eight chips deliver 6.53$\times$ to 7.91$\times$ the
throughput of one, the efficiency rising with sequence length
(Section~\ref{sec:scaling}); against the single-query baseline, which
uses a different input family, the gain is 6.92$\times$.
JAX's automatic partitioner, which is meant to make such manual
mapping unnecessary, leaves the per-chip memory footprint of a
single-query run unchanged from the one-chip case, consistent with
replicating the model rather than splitting it. Our best explanation
is that the model carries no sharding annotations, leaving the
partitioner nothing from which to derive a split, which is the
documented fallback of an annotation-driven partitioner
\citep{xu2021gspmd,shardy2026} and matches the reported effect of
stripping sharding structure from such a system \citep{drjax2024}.
We did not retain the compiler evidence that would confirm it. AlphaFold2's own
ensembling cannot be split this way either, because the public
implementation runs ensemble members in a sequential \texttt{while\_loop}
rather than along an array axis. We build the ensemble outside the
model instead: eight differently seeded featurizations of
one query, one per chip under \texttt{pmap}, with their
\texttt{predicted\_lddt} logits averaged by \texttt{pmean}. AlphaFold2's
source and its internal \texttt{num\_ensemble\,=\,1} path are left
untouched, so this demonstrates that an ensemble can be spread across
the slice; it does not replace AlphaFold2's own ensembling
(Section~\ref{sec:parallelism}).

Prior work has optimized AlphaFold2 training and long-sequence
inference on GPU clusters through model parallelism (FastFold) and through pipeline separation with
compilation reuse (ParaFold) \citep{cheng2024fastfold,zhong2022parafold}.
Both target NVIDIA GPUs; neither reports TPU experiments.
Closer to the hardware, JAXBench includes operators extracted from
AlphaFold2 among its TPU kernel benchmarks \citep{jaxbench2026}, and a recent characterization of
AlphaFold3 profiles its bottlenecks on CPUs and GPUs
\citep{iiswc2025af3}; to our knowledge, neither measures end-to-end
AlphaFold2 inference on TPU.

Our contributions are:
\begin{itemize}
  \item a controlled comparison of AlphaFold2 inference on CPU, GPU and
    a multi-chip TPU slice, using identical code and input, with its
    provenance gaps and cross-session variability reported
    (Sections~\ref{sec:methodology} and~\ref{sec:results-hardware});
  \item a retained trace analysis of the cold-start regime on TPU,
    reporting that about three quarters of the traced \texttt{apply\_fn}
    span is JAX tracing and compilation rather than execution, consistent
    with the compilation cost that ParaFold identified on GPU,
    together with a profiler-overhead correction to our CPU and GPU
    first-call timings, which we could not verify on TPU
    (Section~\ref{sec:profiling});
  \item a measurement showing that vectorization does not raise
    throughput for this model while explicit data parallelism across
    chips does (Section~\ref{sec:parallelism});
  \item an observation that automatic sharding leaves AlphaFold2's
    per-chip footprint unchanged, with a hypothesis for why, and a
    demonstration that an
    ensemble can instead be built outside the model and mapped across
    chips, leaving the model source untouched
    (Section~\ref{sec:parallelism});
  \item a matched 16-point grid over chip count and sequence length,
    with measured eight-over-one-chip speedups of 6.53--7.91$\times$, a
    descriptive fit to it, and its consequences for cost per prediction
    (Sections~\ref{sec:scaling} and~\ref{sec:cost});
  \item a portability finding across model generations: the
    AlphaFold3 release we tested accepts only CPU, GPU and Apple MPS
    backends, not TPU (Section~\ref{sec:alphafold3}).
\end{itemize}

Section~\ref{sec:background} introduces TPUs, the JAX/XLA compilation
model and the computational structure of AlphaFold2.
Section~\ref{sec:methodology} describes the setup, and
Sections~\ref{sec:results-hardware} to~\ref{sec:alphafold3} present the
results. Section~\ref{sec:related-work} places them in context,
Section~\ref{sec:discussion} discusses what they imply for users of
scientific JAX workloads, and Sections~\ref{sec:limitations}
and~\ref{sec:conclusion} state the limitations and conclude.
